%% kapa: exam=1 round=30 year=2019 period=2 form=A subject=accounting subject_ko=회계학 seq=33 number=73
%% source: Q-Net 감정평가사 기출문제 (cst00309) artlSeq=5205205 file=extracted/1008/5205205/2019년 제30회 감정평가사 1차시험 문제지/2019년 제30회 감정평가사 1차시험 문제지/2019년도 제30회 감정평가사 1차 2교시 A형.hwp
%% conversion: kordoc 4.24.51 -> tools/gen1.py (auto); review=proofread-sonnet; flags=none
\begin{question}{73}
(주)감평은 동일한 원재료를 투입하여 제품X, 제품Y, 제품Z를 생산한다. (주)감평은 결합원가를 분리점에서의 상대적 판매가치를 기준으로 결합제품에 배부한다. 결합제품 및 추가가공과 관련된 자료는 다음과 같다. \par
{\small\centering
\begin{tabular}{|c|c|c|c|c|}
\hline
 & \makecell{제품X} & \makecell{제품Y} & \makecell{제품Z} & \makecell{합계} \\
\hline
\makecell{생산량} & \makecell{150단위} & \makecell{200단위} & \makecell{100단위} & \makecell{450단위} \\
\hline
\makecell{결합원가} & \makecell{￦15,000} & \makecell{?} & \makecell{?} & \makecell{?} \\
\hline
\makecell{분리점에서의 단위당 판매가격} & \makecell{￦200} & \makecell{￦100} & \makecell{￦500} &  \\
\hline
\makecell{추가가공원가} & \makecell{￦3,500} & \makecell{￦5,000} & \makecell{￦7,500} & \makecell{￦16,000} \\
\hline
\makecell{추가가공 후 단위당 판매가격} & \makecell{￦220} & \makecell{￦150} & \makecell{￦600} &  \\
\hline
\end{tabular}\par}

(주)감평은 각 제품을 분리점에서 판매할 수도 있고, 분리점 이후에 추가가공을 하여 판매할 수도 있다. (주)감평이 위 결합제품을 전부 판매할 경우, 예상되는 최대 매출총이익은? (단, 결합공정 및 추가가공과정에서 재공품 및 공손은 없다.) \par
\begin{choices}
  \item ￦25,000
  \item ￦57,000
  \item ￦57,500
  \item ￦82,000
  \item ￦120,000
\end{choices}
\end{question}
\answer{3}
